\section{Tetrahedral Relational Frame on $S^2$}
\label{sec:tetrahedral-relational-frame}

\subsection{Regular tetrahedron inscribed in the sphere}
Let
\begin{equation}
T_4 = \{v_0,v_1,v_2,v_3\} \subset S^2
\end{equation}
be a regular tetrahedron on the unit sphere. The defining relations are
\begin{equation}
v_i \cdot v_j =
\begin{cases}
1 & i=j, \\
-\frac{1}{3} & i\neq j,
\end{cases}
\qquad
\sum_{i=0}^3 v_i = 0.
\end{equation}

\subsection{Minimal noncoplanar simplex}
The tetrahedron is the first noncoplanar simplex and therefore the minimal volume-supporting extension beyond surface-only geometry.

\subsection{Relative phase triple}
Assign phases $\phi_i \in U(1)$ to the vertices. After quotienting by global phase, one may choose three independent relative phases,
\begin{equation}
\gamma_1 = \phi_1 - \phi_0,
\qquad
\gamma_2 = \phi_2 - \phi_0,
\qquad
\gamma_3 = \phi_3 - \phi_0.
\end{equation}
This gives the minimal triple of relative phase channels for later closure constructions.

\subsection{Scope restriction}
This section establishes a geometric extension object only. It does not by itself prove a complete physical theory of every appearance of the number $3$.
